% colores-solucion.tex — Colores de los gráficos propios de la solución (1D01,
% examen de España). Se cargan después de tikz-tcee.tex (paleta «dic…»).
% Los 7 primeros son los que añade el cuadernillo del examen (mismo HEX, para
% que lo mismo tenga el mismo color en el examen y en la solución; CRITERIOS §7).
\definecolor{dicbienes}{HTML}{2C6E9B}
\definecolor{dicturismo}{HTML}{8CC3E6}
\definecolor{dicrentas}{HTML}{A6A6A6}
\definecolor{dicsaldo}{HTML}{1A1A1A}
\definecolor{dicpiin}{HTML}{8C510A}
\definecolor{dicfiscal}{HTML}{4A5A6A}
\definecolor{dicvivienda}{HTML}{B07AA1}
% Nuevos en la solución:
% sector privado (gama cálida, la del gasto privado de la guía)
\definecolor{dicprivado}{HTML}{C76A1F}
% efecto bola de nieve (azul medio, distinto del pizarra de lo fiscal)
\definecolor{dicbola}{HTML}{4F86C6}
% ajuste déficit-deuda y otras variaciones (residuos): gris claro con trama
\definecolor{dicresiduo}{HTML}{B4B4B4}
% previsión del PIB (el morado del PIB, aclarado)
\colorlet{dicpibprev}{dicpib!35!white}
\pgfplotsset{
  dicsol/.style={width=13cm, height=6.4cm, scale only axis=false,
    axis lines*=left, tick align=outside, clip=false,
    ymajorgrids, grid style={tceegris!25},
    ticklabel style={font=\footnotesize}, label style={font=\footnotesize},
    xticklabel style={/pgf/number format/1000 sep=},
    yticklabel style={/pgf/number format/.cd, use comma, 1000 sep={}},
    legend style={at={(0.5,-0.15)}, anchor=north, draw=none, font=\footnotesize,
      /tikz/every even column/.append style={column sep=8pt}},
    legend cell align=left},
  leyendalinea/.style={legend image code/.code={\draw[##1] (0cm,0cm) -- (0.55cm,0cm);}},
  % leyenda de una serie de solo puntos (en un gráfico con area legend)
  leyendamarca/.style={legend image code/.code={\draw[mark repeat=2,mark phase=2,##1] plot coordinates {(0cm,0cm) (0.3cm,0cm) (0.6cm,0cm)};}},
}
